\begin{figure}[htbp]
\centering
\begin{tikzpicture}
\begin{axis}[
    width=0.96\textwidth,
    height=6.2cm,
    grid=both,
    grid style={dotted, gray!40},
    xlabel={Messzeit $t$ [Sekunden]},
    ylabel={Fehlerrate [\%]},
    xmin=0, xmax=39,
    ymin=-2, ymax=55,
    legend pos=north west,
    legend style={font=\sffamily\scriptsize, fill=white, fill opacity=0.9},
    label style={font=\sffamily\footnotesize},
    tick label style={font=\sffamily\scriptsize},
    title style={font=\sffamily\bfseries\footnotesize},
]

% Vertikale Markierung fuer Störungsinduktion
\draw[red!80!black, dashed, line width=1pt] (axis cs:15,-2) -- (axis cs:15,55)
    node[pos=0.88, right, font=\sffamily\scriptsize\bfseries, text=red!80!black] {Störungsinduktion ($t = 15\text{ s}$)};

% Basiskonfiguration (1 Pod)
\addplot[
    color=red!75!black,
    mark=square*,
    mark size=1.4pt,
    line width=1.0pt,
    dashed
] table[x=second, y=error_rate, col sep=comma] {daten/exp1-timeseries-1pod.csv};
\addlegendentry{1 Pod (Basiskonfiguration)}

% Verbesserte Konfiguration (2 Pods)
\addplot[
    color=blue!70!black,
    mark=*,
    mark size=1.4pt,
    line width=1.0pt
] table[x=second, y=error_rate, col sep=comma] {daten/exp1-timeseries-2pods.csv};
\addlegendentry{2 Pods (Verbesserte Konfiguration)}

\end{axis}
\end{tikzpicture}
\caption{Zeitlicher Verlauf der sekündlichen Fehlerrate bei forciertem Prozessabsturz ($t=15\text{ s}$) unter Last (40 Anfragen/s), eigene Messung.}
\label{fig:exp1-fehlerrate-verlauf}
\end{figure}
